\subsection{LUSIN SPACES AND BOREL SETS}

THEOREM 3. Let X be a Lusin space. Then a subspace of X is a Lusin space if and only if it is a Borel set.

This theorem is a consequence of the following two lemmas :

LEMMA 6. In a Lusin space X, every Borel set is a Lusin subspace of X.

Let \(\mathfrak{X}\) be the set of all subsets A of X such that both A and \(\complement A\) are Lusin subspaces of X. Since every closed set and every open set in X is a Lusin subspace of X (no. 4), \(\mathfrak{X}\) contains all closed subsets of X, and the lemma will therefore be proved if we can show that \(\mathfrak{X}\) is a \(\sigma\) -algebra on X. For this it is enough to show that if \((A_n)\) is a sequence of sets of \(\mathfrak{X}\) , then \(\bigcap_{n} A_n\) and \(\bigcup_{n} A_n\) are Lusin subspaces of X. Now we have seen in no. 4 that every countable intersection of Lusin subspaces is a Lusin subspace. On the other hand, if \(B_n\) is the intersection of \(A_n\) and \(\bigcap_{i < n} \complement A_i\) , it follows from the hypothesis and the preceding remark that \(B_n\) is a Lusin subspace; and since \(\bigcup_{n} A_n = \bigcup_{n} B_n\) , the subspace \(\bigcup_{n} A_n\) is a Lusin subspace by Lemma 2 of no. 4.

LEMMA 7. Every Lusin subspace A of a metrizable space X is a Borel set in X.

By Proposition 13 of no. 5 there exist a sieve C and a continuous bijection \(f\) of \(\mathbf{L}(\mathbf{C})\) onto A. With the notation of Lemma 5 of no. 6, for each integer \(n\) and each \(c \in \mathbf{C}_n\) , let \(g_n(c)\) denote the subspace \(f(q_n(c))\) of X; it is a Lusin subspace and therefore a Souslin subspace of X. As \(c\) runs through \(\mathbf{C}_n\) , the sets \(g_n(c)\) are pairwise disjoint, because \(f\) is bijective; hence, by Theorem 2 of no. 6, there is a family \(c \to g_n'(c) (c \in \mathbf{C}_n)\) of Borel sets in X, pairwise disjoint and such that \(g_n'(c) \supset g_n(c)\) for all \(c \in \mathbf{C}_n\) . Replacing \(g_n'(c)\) by its intersection with the closure \(\overline{g_n(c)}\) of \(g_n(c)\) in X if necessary, we may suppose that \(g_n'(c) \subset \overline{g_n(c)}\) . Let \(c_{n-1}, c_{n-2}, \ldots, c_0\) denote the images of \(c\) in \(\mathbf{C}_{n-1}, \mathbf{C}_{n-2}, \ldots, \mathbf{C}_0\) , under the surjections

\[
p _ {n - 1, n} = p _ {n - 1}, p _ {n - 2, n} = p _ {n - 2} \circ p _ {n - 1}, \dots , p _ {0 n} = p _ {0} \circ p _ {1} \circ \dots \circ p _ {n - 1}
\]

respectively; and let \(h_n(c)\) denote the intersection of the sets

\[
g _ {n} ^ {\prime} (c), g _ {n - 1} ^ {\prime} \left(c _ {n - 1}\right), \dots , g _ {o} ^ {\prime} \left(c _ {0}\right).
\]

Since \(q_{i}(c_{i}) \supset q_{n}(c)\) for \(0 \leqslant i \leqslant n - 1\) , \(h_{n}(c)\) contains \(g_{n}(c)\) ; it is clear also that \(h_{n}(c)\) is a Borel set and is contained in \(g_{n}(c)\) , and that as \(c\) runs through \(\mathbf{C}_{n}\) , the \(h_{n}(c)\) are mutually disjoint sets; finally, by construction, for each \(c' \in \mathbf{C}_{n+1}\) we have \(h_{n+1}(c') \subset h_n(p_n(c'))\) . Let then \(\mathbf{B}_n\) be the union of the sets \(h_n(c)\) as \(c\) runs through \(\mathbf{C}_n\) ; \(\mathbf{B}_n\) is a Borel set, and \(\mathbf{B}_{n+1} \subset \mathbf{B}_n\) ; also \(\mathbf{B}_n\) contains the union of the sets \(g_n(c) (c \in \mathbf{C}_n)\) , which is A. Let B be the intersection of the decreasing sequence of sets \(\mathbf{B}_n\) ; B is a Borel set and contains A. We shall show that B = A, and this will complete the proof.

Let \(x\) be a point of \(\mathbf{B}\) . Then, for each integer \(n\) , there exists a unique \(c \in \mathbf{C}_n\) such that \(x \in h_n(c)\) ; let us denote this \(c\) by \(c_n(x)\) . The sequence \((c_n(x))_{n \geqslant 0}\) belongs to \(\mathbf{L}(\mathbf{C})\) . The decreasing sequence \((g_n(c_n(x)))\) converges by definition to a point \(a \in \mathbf{A}\) ; the sequence of closures of these sets also converges to \(a\) in \(\mathbf{X}\) , hence a fortiori so does the sequence \((h_n(c_n(x)))\) . Now \(x\) belongs to all the sets \(h_n(c_n(x))\) , therefore \(x = a \in \mathbf{A}\) . Lemma 7 is thus proved, and with it Theorem 3.

COROLLARY. If \(f\) is a continuous injective mapping of a Lusin space (or, in particular, a Polish space) \(\mathbf{X}\) into a metrizable space \(\mathbf{Y}\) , then \(f(\mathbf{X})\) is a Borel set in \(\mathbf{Y}\) .

